\documentclass[11pt]{article}

\usepackage[margin=1in]{geometry}
\usepackage{amsmath}
\usepackage{booktabs}
\usepackage{hyperref}
\usepackage{array}

\title{Statistical Arbitrage in Crypto Pair Spreads\\
\large Working Draft and Paper Outline, Not a Final Manuscript}
\author{Jack Lutz, Sammy Adham, and Frank Kronewitter}
\date{\today}

\begin{document}

\maketitle

\begin{center}
\fbox{\parbox{0.92\linewidth}{
\textbf{Draft status.} This document is a working paper outline prepared from
the current project results. It is not a final paper, not a polished
submission, and not a claim that all analyses are complete. Several claims
below are intentionally phrased as provisional because the project still needs
advisor review, cross-exchange validation, microstructure-aware costs, and
pre-registered out-of-sample evaluation.
}}
\end{center}

\begin{abstract}
This draft summarizes a Math 199 research project on statistical arbitrage in
hourly Binance.US USDT crypto pair spreads. The central methodological result
is that static ordinary least squares hedge ratios fail to produce stable
out-of-sample spreads for most pairs, while a train-only maximum-likelihood
Kalman dynamic hedge ratio produces out-of-sample stationary residuals for
nearly the entire screened liquidity universe. Downstream trading results are
more cautious: machine learning models and z-score rules identify pre-cost
mean-reversion signals, but realistic transaction costs and selection-bias
corrections sharply weaken the profitability claim. At this draft stage, the
defensible conclusion is not that the strategy is tradable, but that dynamic
hedge estimation reveals strong non-stationary relative-price structure whose
economic value remains unresolved.
\end{abstract}

\section*{How to Read This Draft}

This file is an outline-style manuscript. It is meant to collect the current
results in paper form and to make the argument easy to revise. Sections marked
as provisional should be checked before any final submission. Numerical values
come from the repository results documents and scripts as of the current
project state, especially \texttt{docs/RESULTS.md} and
\texttt{docs/METHODS.md}.

\section{Research Question and Motivation}

This project asks whether cryptocurrency prices exhibit stable, quantifiable
relative-value relationships, and whether deviations from those relationships
tend to mean-revert. The research follows the broad pairs-trading framework of
Gatev, Goetzmann, and Rouwenhorst (2006), adapted to a continuously traded
crypto market with a changing universe of listed assets.

The project has two distinct goals. The first is statistical: determine whether
coin-like relationships exist among liquid crypto assets after avoiding obvious
look-ahead bias. The second is economic: test whether any detected
mean-reversion signal remains useful after walk-forward validation,
transaction costs, and selection-bias correction. The current evidence supports
the statistical goal much more strongly than the economic one.

\section{Data and Universe Construction}

The data set contains hourly OHLCV observations from the Binance.US REST API
for USDT spot pairs. Each symbol is stored locally as a parquet file. The
current snapshot ends on 2026-01-22. The primary analysis starts at
\(t_0=\) 2024-01-01, uses a liquidity-filtered universe, and evaluates pairs
with walk-forward train/test splits.

The investible universe is constructed as of each analysis time. A symbol is
eligible only if its first observation is at least the required minimum-history
window before \(t_0\), and its last observation reaches the start of the
analysis period. This first-observation rule was added after an audit found
that post-\(t_0\) listings could otherwise leak into the historical universe.

Survivorship remains a limitation. The local store contains currently listed
symbols, and public Binance.US delisting information implies that roughly
3--6 liquidity-filtered candidates may be missing during the analysis window.
The estimated survivorship-bias bound is about 6.7\% as an upper bound and
about 5\% as a realistic estimate. This does not invalidate the current
results, but it should be treated as a material caveat in the final paper.

\section{Methods}

\subsection{Static and Dynamic Spreads}

The baseline spread uses a static OLS hedge ratio:
\[
s_t = \log(P_{a,t}) - \alpha - \beta \log(P_{b,t}),
\]
where \(\alpha\) and \(\beta\) are fit on the training window.

The main methodological alternative is a two-state Kalman random-walk hedge
model:
\[
\begin{bmatrix}\alpha_t\\ \beta_t\end{bmatrix}
=
\begin{bmatrix}\alpha_{t-1}\\ \beta_{t-1}\end{bmatrix}
+ w_t,
\qquad
\log(P_{a,t}) =
\begin{bmatrix}1 & \log(P_{b,t})\end{bmatrix}
\begin{bmatrix}\alpha_t\\ \beta_t\end{bmatrix}
+ v_t.
\]
The process and observation noise parameters are fit by maximum likelihood on
training data only. For test data, parameters are held fixed and the filter is
rolled forward from the final training state. This distinction is important:
the reported Kalman result is not an in-sample whitening artifact.

\subsection{Prediction Models}

The project compares simple baselines, a z-score rule, gradient boosting, an
LSTM, and transformer variants. Models use backward-looking spread, volume,
correlation, volatility, and market-context features. Labels are three-class:
revert, persist, or diverge over a 24-hour horizon. The headline walk-forward
run uses 27 splits with 90 days of training, 30 days of testing, and a 30-day
step.

\subsection{Backtesting and Statistical Corrections}

Trading evaluation uses a state-machine entry/exit rule: open when the model
predicts reversion and \(|z|\geq 2\), close when \(|z|\leq 0.5\), the model
predicts divergence, or the spread changes sign. Transaction-cost experiments
include zero cost, 5 basis points round-trip per leg, and 15 basis points
round-trip per leg.

The project also reports Newey-West HAC corrections because hourly labels have
overlapping 24-hour targets, and deflated Sharpe ratios because many model and
strategy configurations were tried. These corrections are central to the
current cautious interpretation.

\section{Current Results}

\subsection{Result 1: Static OLS Fails Out-of-Sample}

Static hedge ratios do not produce robust out-of-sample stationary residuals.
Among the top 10 selected pairs, only 2 maintain out-of-sample ADF
\(p<0.05\). Across the full 1,219-pair liquidity-screened universe, static OLS
passes \(p<0.05\) for only 66 pairs, or 5.4\%.

\begin{table}[h]
\centering
\caption{Out-of-sample stationarity screen, static OLS versus Kalman dynamic hedge.}
\begin{tabular}{lrr}
\toprule
ADF threshold & Static OLS passes & Kalman passes \\
\midrule
\(p<0.001\) & 0 (0.0\%) & 1,174 (96.3\%) \\
\(p<0.005\) & 1 (0.1\%) & 1,188 (97.5\%) \\
\(p<0.010\) & 10 (0.8\%) & 1,198 (98.3\%) \\
\(p<0.050\) & 66 (5.4\%) & 1,215 (99.7\%) \\
\(p<0.100\) & 142 (11.6\%) & 1,217 (99.8\%) \\
\bottomrule
\end{tabular}
\end{table}

\subsection{Result 2: Kalman Dynamic Hedge Ratios Recover Stationarity}

The strongest current finding is that dynamic hedge ratios recover
stationary-looking residuals for nearly all liquid pairs under an honest
train/test protocol. On the full 1,219-pair screen, 99.7\% of pairs pass
\(p<0.05\), and 96.3\% pass \(p<0.001\). The selected top-10 pairs all pass
\(p<0.001\) under the Kalman model.

This should be the leading result in the final paper unless advisor review
identifies a flaw in the protocol. The interpretation is that crypto pair
relationships are real but non-stationary: a fixed hedge ratio is usually too
rigid, while a slowly drifting hedge ratio captures the relative-price
structure.

\subsection{Result 3: ML Helps, but the Main Signal Is Still Spread Z-Score}

In the post-audit 27-split walk-forward run, the LSTM has the strongest
\texttt{pnl\_mean\_to\_std}, followed by the booster and z-score rule. The
small parameter-matched transformer is not competitive in this run, although a
larger transformer later performs much better on a related pre-audit data set.

\begin{table}[h]
\centering
\caption{Post-audit walk-forward model comparison. Values are current draft results.}
\begin{tabular}{lrrr}
\toprule
Model & \texttt{pnl\_mean\_to\_std} & Win rate & Trades/split \\
\midrule
LSTM & 0.376 & 0.795 & 2,083 \\
Booster & 0.356 & 0.741 & 2,355 \\
Z-score & 0.282 & 0.963 & 588 \\
Majority & 0.021 & 0.028 & 185 \\
Small transformer & 0.018 & 0.437 & 691 \\
Random & 0.001 & 0.502 & 2,689 \\
\bottomrule
\end{tabular}
\end{table}

A feature ablation suggests that the learned models are mostly refining a
classical z-score signal. Dropping \texttt{latest\_spread\_z} reduces accuracy
by 14 percentage points, macro F1 by 19 percentage points, and PnL by roughly
79\%. Other features have much smaller effects in the current single-split
ablation. This is an important draft conclusion: the ML models do not yet
demonstrate a qualitatively new source of alpha beyond the spread deviation
itself.

\subsection{Result 4: Cost and Selection Bias Substantially Weaken Trading Claims}

The strategy looks strong before costs, with annualized Sharpe ratios around
8 under the state-machine backtest. That headline does not survive careful
interpretation. At 5 bps round-trip per leg, Sharpe falls to roughly 0.8--0.9.
At 15 bps round-trip per leg, the strategy is deeply negative.

\begin{table}[h]
\centering
\caption{Cost sensitivity under the state-machine backtest.}
\begin{tabular}{lrrr}
\toprule
Cost per leg & LSTM Sharpe & Z-score Sharpe & Booster Sharpe \\
\midrule
0 bps & 8.08 & 8.03 & 8.27 \\
5 bps & 0.80 & 0.83 & 0.90 \\
15 bps & -14.04 & -13.88 & -14.17 \\
\bottomrule
\end{tabular}
\end{table}

Deflated Sharpe correction further weakens the profitability claim because the
project tried many strategy configurations. At a modest trial count, the
pre-cost result may remain reportable. At a conservative count around 50
trials, even the pre-cost result is not clearly distinguishable from the
chance-best result. The 5 bps result does not survive deflated-Sharpe
correction at any plausible trial count. The current paper should therefore
retract any break-even cost claim.

\section{Negative Results}

The HMM regime filter is a robust negative result. In post-audit runs, the
filter suppresses 40--55\% of bars and reduces total PnL across models rather
than isolating productive mean-reversion regimes. This remains true under
2-state and 3-state variants and across static and Kalman pipelines.

Single-split evaluation is also unreliable. An earlier single chronological
split overstated the transformer result. The 27-split walk-forward evaluation
reduced the small transformer's apparent performance until its confidence
interval included zero. This supports treating walk-forward evaluation with
uncertainty intervals as the minimum reporting standard for this project.

\section{Audit Findings}

The project found and fixed several methodological issues during review:
\begin{itemize}
\item post-\(t_0\) listings could enter the historical universe.
\item an in-sample Kalman ADF check was circular and was replaced by a
train/test protocol.
\item per-pair label thresholds could use information beyond the relevant
training boundary.
\item the Kalman spread construction had a train/test parameter discontinuity.
\item the backtest Sharpe denominator could be distorted by changing active
capital.
\item \texttt{time\_since\_zero\_crossing} was confirmed as a soft leakage
feature and removed from the post-audit interpretation.
\end{itemize}

The repository currently includes leakage and backtester sanity tests. The
draft should report these as evidence of methodological cleanup, not as proof
that all remaining bias has been eliminated.

\section{Limitations}

This draft should retain the following limitations prominently:
\begin{itemize}
\item Results use one exchange, Binance.US.
\item Results use one quote currency, USDT.
\item Results use one bar interval, hourly.
\item Survivorship bias remains because delisted assets are missing from the
local store.
\item The backtester uses flat slippage and does not model market impact.
\item Spot shorting and funding constraints are not fully modeled.
\item All main results use \(t_0=\) 2024-01-01. Start-date sensitivity remains
pending.
\item Deep-learning results are mostly single-seed.
\item No pre-registered final strategy has been evaluated exactly once.
\end{itemize}

\section{Provisional Discussion}

The current results suggest a distinction between statistical structure and
tradable alpha. Dynamically hedged crypto pair spreads appear highly
mean-reverting in a statistical sense, especially compared with static OLS
spreads. However, converting that structure into a realistic trading strategy
is much harder. Transaction costs, selection bias, and market microstructure
could plausibly absorb the apparent edge.

This is still a useful result. A final paper can argue that the appropriate
model for crypto pair relationships is time-varying rather than static, while
also showing that statistical mean reversion is not sufficient evidence of
profitable arbitrage. That framing is more defensible than claiming a finished
profitable strategy.

\section{Work Needed Before Final Submission}

Before this becomes a final manuscript, the project should ideally complete:
\begin{enumerate}
\item advisor review of the Kalman OOS protocol.
\item a rerun of downstream pair selection using Kalman-screened pairs rather
than correlation fallback.
\item cross-exchange validation on Coinbase or Kraken.
\item L2-data ingestion and market-impact modeling.
\item multi-seed deep-learning runs.
\item start-date sensitivity tests.
\item one pre-registered model configuration evaluated once.
\end{enumerate}

\section{Draft Conclusion}

At this stage, the project supports a cautious conclusion. Static OLS hedge
ratios are inadequate for hourly Binance.US crypto pair spreads, while Kalman
dynamic hedge ratios recover strong out-of-sample stationary residuals across
nearly the entire liquid universe. Predictive models can exploit this signal
before costs, but the resulting trading performance does not yet survive
realistic cost assumptions and selection-bias correction strongly enough to
claim deployable arbitrage. The final paper should emphasize the dynamic
hedging result and present the trading results as preliminary, conditional,
and limited by current data and execution assumptions.

\begin{thebibliography}{9}

\bibitem{gatev2006}
Gatev, E., Goetzmann, W. N., and Rouwenhorst, K. G. (2006).
\textit{Pairs Trading: Performance of a Relative-Value Arbitrage Rule}.
Review of Financial Studies, 19(3), 797--827.

\bibitem{bailey2014}
Bailey, D. H., and Lopez de Prado, M. (2014).
\textit{The Deflated Sharpe Ratio: Correcting for Selection Bias,
Backtest Overfitting, and Non-Normality}.

\end{thebibliography}

\end{document}
